\documentclass[12pt]{article}
\usepackage[utf8]{inputenc}
\usepackage[T1]{fontenc}
\usepackage{hyperref}
\title{Cisalpine Gaulish Language Resources for K-12 Curriculum Development}
\author{Compiled by Claude Sonnet 5.5 (claude-sonnet-5-5) for Andr\'as Kornai}
\date{}
\begin{document}
\maketitle

\section{Introduction}
Cisalpine Gaulish (ISO 639-3 xcg) is an extinct Celtic language of northern Italy. It is known only from a small corpus of short inscriptions. One source counts about 20 Cisalpine Gaulish inscriptions, and only five of them are longer than one or two words. Scholars disagree about whether it is a dialect of Gaulish or a continuation of Lepontic. It has no native speakers, and I found no school programmes, textbooks or revitalization materials. A K-12 curriculum would have to be built from the scholarly and epigraphic literature.

\section{Primary Educational Resources}
I found no primers, textbooks or classroom materials for Cisalpine Gaulish.

\section{Linguistic and Reference Materials}
\begin{itemize}
\item Wikipedia, ``Cisalpine Gaulish'': general overview and bibliography. \url{https://en.wikipedia.org/wiki/Cisalpine_Gaulish}
\item Wikipedia, ``Lepontic language'': the related and sometimes overlapping Celtic language of the same area. \url{https://en.wikipedia.org/wiki/Lepontic_language}
\item Lexicon Leponticum (LexLep): a digital edition of the Lepontic and Cisalpine Gaulish inscriptions with an etymological dictionary. The editors include Corinna Salomon, David Stifter and Martin Braun. The Digital Classicist wiki describes it. \url{https://wiki.digitalclassicist.org/Lexicon_Leponticum}. The search summary named the project's own site as lexlep.univie.ac.at, but that address did not appear as a result link, so it is not listed as a resource entry.
\item David Stifter, ``Cisalpine Celtic'' (Palaeohispanica 20, 2020): a survey of the corpus of about 430 short texts in Lepontic and Cisalpine Gaulish. It is freely available through Maynooth University's repository. \url{https://mural.maynoothuniversity.ie/id/eprint/14948/1/DS-Cisalpine-Celtic-2020.pdf}
\item The same article is also on Academia.edu. \url{https://www.academia.edu/44688418/Cisalpine_Celtic}
\item Stifter's \emph{Cisalpine Celtic: Language, Writing, Epigraphy} (Aelaw Booklet 8) and \emph{Old Celtic Languages} (2008) are cited in the Wikipedia bibliography. I saw no direct URL for either.
\end{itemize}

\section{Religious and Community Resources}
I found no religious or community resources for Cisalpine Gaulish.

\section{Digital and Audio Resources}
I found no audio or video materials. The only digital resources found are the LexLep database and the open-access article above.

\section{Curriculum Development Context}
There are no speakers and no existing programmes. Teachers could use the epigraphic corpus to build short units on the history of northern Italy, writing systems (the inscriptions use the North Italic alphabet), and Celtic languages. Any activity should say plainly that the evidence is a small set of short inscriptions and that the classification of the language is disputed.

\end{document}